% formulasheet.tex -- Module 7: Formula sheet explained.
% The first table of each chapter = the formula cards of that chapter (same form, same
% badge: default = on the sheet, [notonsheet], [combine=...]). "Also used" = formulas used in
% the chapter's text, worked examples and boxes without a card of their own.
\unnumberedchapter[\faIcon{square-root-alt}]{Formula sheet explained}\label{fs:start}

This chapter lists every formula used in Module~7, chapter by chapter. Each row says what the
formula is for and when to use it. The last column tells you about the official HSC formula
sheet, which you are given in the exam:
\begin{itemize}
  \item \textcolor{hscformula}{\footnotesize\bfseries\faCheck\ Yes}: it is printed on the
        sheet, in exactly this form;
  \item \textcolor{hscmistake}{\footnotesize\bfseries\faTimes\ Learn it}: it is not on the
        sheet, so learn it or know how to derive it (the formula card shows how);
  \item \textcolor{hsctiptext}{\footnotesize\bfseries\faLink\ Combine}: it is not printed in
        this form, but you can get it by joining formulas (or data) from the sheet.
\end{itemize}
The rearranged forms of each formula are on its formula card, so they are not listed again
here. The sheet prints $v = f\lambda$, not $c = f\lambda$: for light in a vacuum, put
$v = c$.

\section*{Chapter 1: The electromagnetic spectrum}

\begin{formulatable}
  \frow[notonsheet]{$c = \dfrac{1}{\sqrt{\mu_0\varepsilon_0}}$}{Maxwell's predicted speed of EM
    waves, from two data-sheet constants. Not for light in glass or water.}
  \frow[combine={$v = f\lambda$ with $v = c$}]{$c = f\lambda$}{any EM wave in a vacuum or
    air (combine $v = f\lambda$ with $v = c$). Change \si{nm} to \si{m} and \si{GHz} to
    \si{Hz} first.}
  \frow[notonsheet]{$c = \dfrac{d}{t}$}{timing light over a known distance (Rømer, Fizeau,
    radar, laser ranging). A signal there and back travels $d = 2D$.}
  \frow{$\lambda_\mathrm{max} = \dfrac{b}{T}$}{Wien's law: the peak of a star's continuous
    spectrum gives its surface temperature. $T$ in kelvin; not for one spectral line.}
\end{formulatable}

\subsection*{Also used in Chapter 1}
\begin{formulatable}
  \frow{$v = f\lambda$}{any wave: speed $=$ frequency $\times$ wavelength.}
  \frow{$f = \dfrac{1}{T}$}{frequency from the period of one oscillation.}
  \frow[notonsheet]{$t = \dfrac{1}{2Nf}$}{Fizeau's wheel ($N$ teeth, $f$ turns each second):
    the time for a gap to turn to the next tooth. Derived on the $c = d/t$ card.}
  \frow[notonsheet]{$\Delta\lambda = \dfrac{\lambda v}{c}$}{the size of the Doppler shift of a
    star's line, with $v$ its speed along the line of sight (only for $v$ much smaller than $c$;
    HSC questions usually ask about shifts in words). Never use the sheet's Doppler formula for
    light: it is for sound.}
  \frow[notonsheet]{$T(\unit{K}) = T(\unit{\degreeCelsius}) + 273$}{changes a temperature in
    \si{\degreeCelsius} into kelvin before you use Wien's law.}
  \frow{$E = hf$}{the photon from one energy-level gap gives one spectral line (card in
    \cref{ch3:chap}).}
\end{formulatable}

\section*{Chapter 2: Light as a wave}

\begin{formulatable}
  \frow{$n_\mathrm{x} = \dfrac{c}{v_\mathrm{x}}$}{the speed of light in water or glass from
    the refractive index. $n \ge 1$, so $v_\mathrm{x}$ is never more than $c$.}
  \frow{$d\sin\theta = m\lambda$}{bright fringes of a double slit or a grating. $d$ = slit
    separation; $m$ = order (0, 1, 2, \ldots).}
  \frow[notonsheet]{$d\sin\theta = \left(m - \tfrac{1}{2}\right)\lambda$\newline
    {\footnotesize $m = 1, 2, 3, \ldots$}}{dark fringes of a double slit. Never put $m = 0$
    in this formula.}
  \frow[notonsheet]{$\delta = \left|S_2P - S_1P\right|$}{path difference from the two distances
    to a point P (two in-phase sources, near or far). A whole number of wavelengths:
    constructive. A whole number plus a half: destructive.}
  \frow[notonsheet]{$y = L\tan\theta$}{changes a fringe angle into a distance on the screen,
    or back. Works at any angle. Measure $y$ from the centre.}
  \frow[notonsheet]{$\Delta y = \dfrac{\lambda L}{d}$}{the gap between neighbouring fringes of
    a double slit, at small angles (below about \ang{10}). Not for most gratings.}
  \frow[notonsheet]{$d = \dfrac{1}{N}$}{grating spacing from $N$ lines per metre (lines per
    \si{mm} $\times 1000$).}
  \frow[notonsheet]{$I = \tfrac{1}{2}I_0$}{unpolarised light through the first polariser, at
    any angle.}
  \frow{$I = I_\mathrm{max}\cos^2\theta$}{Malus' law: light that is already polarised passes
    a polariser. $\theta$ is measured from the filter before it.}
\end{formulatable}

\subsection*{Also used in Chapter 2}
\begin{formulatable}
  \frow{$n_1\sin\theta_1 = n_2\sin\theta_2$}{Snell's law (Year 11): the angle of a refracted
    ray.}
  \frow{$v = f\lambda$}{the wavelength of sound or light from its speed and frequency.}
  \frow[notonsheet]{$v_1\sin\theta_1 = v_2\sin\theta_2$}{Newton's particle model of
    refraction (the velocity along the surface stays the same). It wrongly predicts faster
    light in water.}
  \frow[combine={$s = ut + \frac{1}{2}at^2$ with $a = 0$}]{$t = \dfrac{s}{v}$}{the time for
    light to cross a distance (combine $s = ut + \frac{1}{2}at^2$ with $a = 0$).}
  \frow[notonsheet]{$\delta = d\sin\theta$}{the path difference for two slits: bright where
    $\delta = m\lambda$, dark where $\delta = \left(m - \frac{1}{2}\right)\lambda$.}
  \frow[notonsheet]{highest $m$ $=$ $\dfrac{d}{\lambda}$\newline{\footnotesize rounded
    down}}{the highest order a grating shows ($\sin\theta$ cannot be more than 1). If
    $d/\lambda$ is a whole number, that order is at \ang{90}: use one less.}
\end{formulatable}

\section*{Chapter 3: Light as a particle}

\begin{formulatable}
  \frow{$\lambda_\mathrm{max} = \dfrac{b}{T}$}{Wien's law for black bodies: the hotter the
    object, the shorter its peak wavelength. $T$ in kelvin.}
  \frow{$E = hf$}{the energy of one quantum (photon). A beam's total energy is (number of
    photons) $\times\, hf$.}
  \frow[combine={$W = qV$}]{$K_\mathrm{max} = qV_0$}{changes a stopping voltage into the
    maximum kinetic energy, or back (from $W = qV$). $K_\mathrm{max}$ in \si{eV} has the same
    number as $V_0$ in volts.}
  \frow{$K_\mathrm{max} = hf - \phi$}{Einstein's photoelectric equation, when $hf \ge \phi$.
    Use all joules or all \si{eV}.}
  \frow[notonsheet]{$\phi = hf_0 \qquad \lambda_0 = \dfrac{hc}{\phi}$}{the threshold:
    electrons leave only if $f \ge f_0$, which means $\lambda \le \lambda_0$.}
  \frow[notonsheet]{$\dfrac{N}{t} = \dfrac{P}{hf} = \dfrac{P\lambda}{hc}$}{the number of photons
    each second from a one-colour lamp or laser of power $P$. Change \si{mW} to \si{W}.}
  \frow[notonsheet]{$V_0 = \dfrac{h}{q}f - \dfrac{\phi}{q}$}{a graph of $V_0$ against $f$:
    gradient $= h/q$, $f$-intercept $= f_0$, $V_0$-intercept $= -\phi/q$.}
\end{formulatable}

\subsection*{Also used in Chapter 3}
\begin{formulatable}
  \frow[combine={$E = hf$ and $v = f\lambda$}]{$E = \dfrac{hc}{\lambda}$}{photon energy from
    the wavelength (shown on the $E = hf$ card; combine $E = hf$ and $v = f\lambda$). Change
    \si{nm} to \si{m}; the answer is in joules.}
  \frow[combine={$v = f\lambda$ with $v = c$}]{$c = f\lambda$}{frequency from wavelength for
    light (card in \cref{ch1:chap}).}
  \frow{$W = qV$}{the work done on a charge $q$ by a voltage $V$; it defines the electron
    volt.}
  \frow{$P = \dfrac{\Delta E}{\Delta t}$}{power: the energy given out each second.}
  \frow{$K = \dfrac{1}{2}mv^2$}{the speed of the fastest photoelectrons (put $K_\mathrm{max}$
    in joules).}
  \frow[notonsheet]{$T(\unit{K}) = T(\unit{\degreeCelsius}) + 273$}{changes a temperature in
    \si{\degreeCelsius} into kelvin before you use Wien's law.}
  \frow[notonsheet]{$0,\ hf,\ 2hf,\ 3hf, \ldots$}{Planck's idea: an oscillator's energy can
    only change in steps of $hf$ (an idea, not for calculations).}
\end{formulatable}

\section*{Chapter 4: Light and special relativity}

\begin{formulatable}
  \frow{\rule[-5.0ex]{0pt}{0pt}$t = \dfrac{t_0}{\sqrt{\left(1 - \dfrac{v^2}{c^2}\right)}}$}{time
    dilation. $t_0$ = proper time, measured by the clock that is present at both events.
    Check that $t > t_0$.}
  \frow{\rule[-2.6ex]{0pt}{0pt}$l = l_0\sqrt{\left(1 - \dfrac{v^2}{c^2}\right)}$}{length
    contraction along the motion. $l_0$ = proper length, measured where the object is at
    rest. Check that $l < l_0$.}
  \frow{\rule[-5.0ex]{0pt}{0pt}$p_\mathrm{v} = \dfrac{m_0 v}{\sqrt{\left(1 -
    \dfrac{v^2}{c^2}\right)}}$}{momentum at any speed; above about $0.1c$ you must use it. Put
    $v$ in \si{m.s^{-1}}. Not for photons.}
  \frow{$E = mc^2$}{energy from lost mass: $\Delta E = \Delta m\,c^2$ with
    $\Delta m = m_\mathrm{before} - m_\mathrm{after}$ (mass in \si{kg}). Annihilation of an
    electron and a positron: $\Delta m = 2m_\mathrm{e}$, two photons.}
  \frow[combine={$P = \frac{\Delta E}{\Delta t}$ and $E = mc^2$}]{$\dfrac{\Delta m}{\Delta t}
    = \dfrac{P}{c^2}$}{the mass lost each second by the Sun, a star or a power station of
    power $P$ (combine $P = \frac{\Delta E}{\Delta t}$ and $E = mc^2$).}
\end{formulatable}

\subsection*{Also used in Chapter 4}
\begin{formulatable}
  \frow{$P = \dfrac{\Delta E}{\Delta t}$}{power: the energy given out each second (the Sun).}
  \frow[combine={$E = mc^2$ and $u = 931.5$ MeV$/c^2$}]{$E\,(\mathrm{MeV}) = \Delta m\,
    (\mathrm{u}) \times 931.5$}{energy from a mass change given in \si{\amu} (data sheet:
    \qty{1}{\amu} $=$ \qty{931.5}{\mega\eV}$/c^2$). Never multiply a mass in \si{\amu} by
    $c^2$.}
  \frow{$E = hf$}{the energy of each photon from annihilation.}
  \frow{$\Delta\vec{p} = \vec{F}_\mathrm{net}\Delta t$}{a finite force for a finite time gives
    only a finite momentum, so a mass never reaches $c$.}
  \frow{$K = \dfrac{1}{2}mv^2$}{classical kinetic energy, for comparison only: it is wrong
    near $c$.}
  \frow[combine={$s = ut + \frac{1}{2}at^2$ with $a = 0$}]{$s = vt$}{distance at a constant
    speed (muon problems): use dilated time with proper distance, or proper time with
    contracted distance.}
  \frow[notonsheet]{$t_0 = \dfrac{2d}{c}$}{one tick of a light clock (mirrors $d$ apart) in its
    own frame; used to derive time dilation.}
  \frow[notonsheet]{time in years $=$\newline distance in \si{\lightyear} $\div \dfrac{v}{c}$}{a
    quick trip time in the Earth frame (\qty{1}{\lightyear} is about \qty{9.46e15}{m}).}
  \frow[notonsheet]{velocity on the ground $=$\newline frame velocity $+$ velocity in the
    frame}{everyday (Galilean) adding of velocities. Never for light: light always travels at
    $c$.}
  \frow[notonsheet]{$r = \dfrac{mv}{qB} = \dfrac{p}{qB}$}{the radius of a charged particle's
    path in a magnetic field (from $qvB = mv^2/r$). Near $c$, use $p = p_\mathrm{v}$.}
  \frow{$v = \dfrac{2\pi r}{T}$}{the speed of a point on the equator as the Earth spins
    ($r = r_\mathrm{E}$, $T$ = one day in seconds), for clocks flown east and west.}
  \frow[notonsheet]{$t - t_0 \approx t \times \dfrac{v^2}{2c^2}$}{time dilation when $v$ is
    tiny compared with $c$ (aeroplanes, Hafele--Keating), when a calculator shows the root
    as 1.}
  \frow[notonsheet]{$1 + z = \dfrac{\lambda_\mathrm{observed}}{\lambda_\mathrm{emitted}}$}{the
    stretch factor of a distant supernova: its light curve lasts $1 + z$ times longer.}
\end{formulatable}

\section*{Data sheet values used in Module 7}

These values are on the official data sheet, exactly as printed. You do not need to learn
them, but you must know which one to use.

{\small
\renewcommand\arraystretch{1.3}\arrayrulecolor{hscrule}%
\setlength\LTpre{4pt}\setlength\LTpost{4pt}%
\begin{longtable}{@{}L{37mm}@{\hspace{3mm}}C{12mm}@{\hspace{3mm}}L{47mm}@{\hspace{3mm}}L{41mm}@{}}
  \colhead{Quantity} & \colhead{Symbol} & \colhead{Value} & \colhead{Used for}\\\hline
  \endfirsthead
  \colhead{Quantity} & \colhead{Symbol} & \colhead{Value} & \colhead{Used for}\\\hline
  \endhead
  Charge on electron & $q_\mathrm{e}$ & \qty{-1.602e-19}{C} & $K_\mathrm{max} = qV_0$, $W = qV$
    (use the size, without the minus sign)\\\hline
  Mass of electron & $m_\mathrm{e}$ & \qty{9.109e-31}{kg} & photoelectron speed; $p_\mathrm{v}$;
    annihilation\\\hline
  Mass of proton & $m_\mathrm{p}$ & \qty{1.673e-27}{kg} & $p_\mathrm{v}$ and rest energy of a
    proton\\\hline
  Speed of sound in air & & \qty{340}{m.s^{-1}} & sound diffracting through a doorway\\\hline
  Speed of light & $c$ & \qty{3.00e8}{m.s^{-1}} & almost every chapter\\\hline
  Electric permittivity constant & $\varepsilon_0$ & \qty{8.854e-12}{A^2.s^4.kg^{-1}.m^{-3}} &
    $c = 1/\sqrt{\mu_0\varepsilon_0}$\\\hline
  Magnetic permeability constant & $\mu_0$ & $4\pi\times10^{-7}$ \si{N.A^{-2}} &
    $c = 1/\sqrt{\mu_0\varepsilon_0}$\\\hline
  Radius of Earth & $r_\mathrm{E}$ & \qty{6.371e6}{m} & the speed of the ground at the equator
    (clocks flown east and west)\\\hline
  Planck constant & $h$ & \qty{6.626e-34}{J.s} & $E = hf$, $K_\mathrm{max} = hf - \phi$\\\hline
  Atomic mass unit & $u$ & \qty{1.661e-27}{kg}\newline \qty{931.5}{\mega\eV}$/c^2$ & mass
    changes given in \si{\amu}\\\hline
  1 eV & & \qty{1.602e-19}{J} & changing between \si{eV} and \si{J}\\\hline
  Wien's displacement constant & $b$ & \qty{2.898e-3}{m.K} & $\lambda_\mathrm{max} = b/T$\\\hline
\end{longtable}}

\section*{Symbol traps}

Some letters on the sheet mean different things in different formulas.

{\small
\renewcommand\arraystretch{1.3}\arrayrulecolor{hscrule}%
\setlength\LTpre{4pt}\setlength\LTpost{4pt}%
\begin{longtable}{@{}C{18mm}@{\hspace{3mm}}L{60mm}@{\hspace{3mm}}L{64mm}@{}}
  \colhead{Symbol} & \colhead{In Module 7} & \colhead{Elsewhere on the sheet}\\\hline
  \endfirsthead
  \colhead{Symbol} & \colhead{In Module 7} & \colhead{Elsewhere on the sheet}\\\hline
  \endhead
  $T$ & temperature in kelvin in $\lambda_\mathrm{max} = b/T$; but the period (the time for one
    cycle) in $f = 1/T$ and $v = 2\pi r/T$ & the period in $f = 1/T$, $v = 2\pi r/T$ and
    $r^3/T^2 = GM/4\pi^2$ (Module 5)\\\hline
  $N$ & teeth on Fizeau's wheel; lines in each metre of a grating ($d = 1/N$); the number of
    photons in $N/t$ & the number of turns of a coil (Module 6); $N_\mathrm{t}$ and $N_0$,
    numbers of nuclei (Module 8)\\\hline
  $u$ & the atomic mass unit: \qty{1}{\amu} $=$ \qty{1.661e-27}{kg} $=$
    \qty{931.5}{\mega\eV}$/c^2$ & the initial velocity in $s = ut + \frac{1}{2}at^2$ and
    $v = u + at$ (Module 5)\\\hline
  $m$ & the order of a fringe in $d\sin\theta = m\lambda$ & mass (as in $E = mc^2$)\\\hline
  $d$ & slit (or grating) spacing in $d\sin\theta = m\lambda$; in this book also a distance
    ($c = d/t$) and the mirror gap of a light clock & the distance between two plates in
    $E = V/d$ (Module 6)\\\hline
  $I$ & light intensity in Malus' law (\si{W.m^{-2}}) & electric current (Module 6)\\\hline
  $E$ & energy ($E = hf$, $E = mc^2$) & electric field strength (Module 6)\\\hline
  $\lambda$ & wavelength & the decay constant (Module 8)\\\hline
  $\phi$ & work function (lower-case phi) & $\Phi$ (capital) is magnetic flux
    (Module 6)\\\hline
  $c$ & the speed of light & specific heat capacity in $Q = mc\Delta T$ (Year 11)\\\hline
  $l$ & a length along the motion, in length contraction & wire length
    (Module 6)\\\hline
  $t_0$, $l_0$, $m_0$ & proper time, proper length, rest mass & \\\hline
\end{longtable}}
